\Needspace{9\baselineskip}
\originalchapter{量子电路与信息可读出性}
\originalsection{寄存器、门与可逆嵌入}
$n$比特计算基$\ket x$, $x\in\{0,1\}^n$, 态为$\sum_x\alpha_x\ket x$. 有$2^n$个幅度并不意味着能一次读出$2^n$个数; 计算基测量只产生一个$x$, 概率$|\alpha_x|^2$. 算法的任务是让目标函数决定不同路径的相位, 再利用干涉把所需全局性质集中到少量可读结果上.

不可逆经典函数$f$通过$U_f\ket x\ket y=\ket x\ket{y\oplus f(x)}$嵌入可逆操作. 输入$x$保留才能使映射一一, 工作垃圾寄存器若与答案纠缠会抑制后续干涉. 计算$f$、把答案复制到新基态寄存器、再逆转工作电路可清除垃圾; 此“复制”只对正交计算基标签有效, 不构成未知态克隆.

Hadamard $H=(X+Z)/\sqrt2$把$0$变成$(0+1)/\sqrt2$. CNOT $\ket{x,y}\mapsto\ket{x,y\oplus x}$产生纠缠, 其共轭作用$X_c\mapsto X_cX_t,Z_t\mapsto Z_cZ_t,X_t\mapsto X_t,Z_c\mapsto Z_c$将在纠错中描述误差传播. Clifford门保持Pauli群, 可高效模拟; 加非Clifford门如$T=\operatorname{diag}(1,\e^{\ii\pi/4})$才突破该可模拟结构. “生成纠缠”并非计算加速的充分条件.

\originalsection{相位反冲与简单算法}
把函数目标寄存器置$\ket-=(0-1)/\sqrt2$, 则$X^{f(x)}\ket-=(-1)^{f(x)}\ket-$, 黑盒使输入产生相位$(-1)^{f(x)}$. Deutsch--Jozsa承诺$f$恒定或恰半数为1, 初始均匀输入, 黑盒相位后再作$H^{\otimes n}$, 零串振幅为$2^{-n}\sum_x(-1)^{f(x)}$. 恒定时绝对值1, 平衡时0. 加速依赖承诺和查询模型, 不能推广成能高效求任意$f$的全部值.

Simon承诺$f(x)=f(y)$当且仅当$y=x$或$x\oplus s$. 测函数寄存器后输入为$(\ket x+\ket{x\oplus s})/\sqrt2$. Hadamard后$\ket z$振幅含$1+(-1)^{z\cdot s}$, 所以只读到$z\cdot s=0$的线性方程. 反复取得约$n$个独立方程, 用$\mathbb F_2$消元求$s$. 测函数标签的具体值无关, 它只选择一个陪集; 这个陪集结构在隐藏子群算法再次出现.

\originalsection{实现、资源与精度}
电路复杂度需同时计逻辑门数、工作比特、测量和允许误差. 模块化幂运算仍是一段需实现的可逆算术, 黑盒不能凭名称免费获得. 有限普适门集逼近连续旋转要额外编译成本. 实验实现还要求可辨识的两态、初始化、受控耦合、读出与足够长相干时间. 离子阱、原子与光子实现的物理耦合见原子光物理卷, 本卷只讨论相应门模型和误差后果, 不重复量子光场基础.
